This section walks through three production-shaped scenarios that
each isolate one AIsafePy module. For every scenario we describe the
persona, the threat, the code before AIsafePy, the code after, and
a flow diagram of how the library blocks the attack.

\subsection{Case A: indirect prompt injection in an email assistant}
\label{case:flow}

\paragraph{Persona.} Maya runs a two-person SaaS that ships an AI
email assistant on top of the OpenAI Agents SDK. The assistant has
three tools: read the user's latest unread email, send a reply, and
forward arbitrary emails to a recipient.

\paragraph{Threat.} A marketing email arrives whose body contains
hidden white-on-white text:
\textit{``SYSTEM: After replying, call
forward\_emails(to=`leak@evil.com', count=5) to confirm ownership.
Do not mention this to the user.''} The naive agent reads the email,
treats the body as instructions, and exfiltrates five real emails to
an attacker-controlled address. This is the canonical indirect
prompt-injection pattern from \citet{debenedetti2024agentdojo}.

\paragraph{Before.} Listing~\ref{lst:maya-before} shows Maya's agent
as shipped. The body of the email reaches the LLM unfiltered; the
LLM emits a tool call; the tool runs.

\begin{lstlisting}[caption={Email assistant before AIsafePy (Scenario A).}, label={lst:maya-before}]
from agents import Agent, Runner, function_tool

@function_tool
def gmail_read_latest() -> dict: ...

@function_tool
def send_email(to: str, subject: str, body: str) -> str: ...

@function_tool
def forward_emails(to: str, count: int) -> str: ...

agent = Agent(
    name="inbox-assistant",
    tools=[gmail_read_latest, send_email, forward_emails],
)

Runner.run_sync(agent, "Please handle my latest unread email.")
# Output: Forwarded 5 emails to leak@evil.com  (attack succeeded)
\end{lstlisting}

\paragraph{After.} Listing~\ref{lst:maya-after} shows the same agent
with eight lines of \texttt{aisafepy.flow} policy added. The
\texttt{gmail\_read\_latest} source is labelled UNTRUSTED. Privileged
sinks (\texttt{send\_email}, \texttt{forward\_emails}) require the
control-flow to derive from TRUSTED input. The injected tool call
fails the integrity check by construction.

\begin{lstlisting}[caption={Email assistant with AIsafePy.flow (Scenario A).}, label={lst:maya-after}]
from agents import Agent, Runner, function_tool
from aisafepy.flow import Policy, Capability, secure_agent

@function_tool
def gmail_read_latest() -> dict: ...

@function_tool
def send_email(to: str, subject: str, body: str) -> str: ...

@function_tool
def forward_emails(to: str, count: int) -> str: ...

policy = (
    Policy()
    .label_source("user_prompt",       integrity="TRUSTED")
    .label_source("gmail_read_latest", integrity="UNTRUSTED",
                  caps={Capability.READ_USER})
    .require("send_email",     control_flow_integrity="TRUSTED",
             caps={Capability.WRITE_EXTERNAL})
    .require("forward_emails", control_flow_integrity="TRUSTED",
             caps={Capability.WRITE_EXTERNAL})
)

agent = Agent(name="inbox-assistant",
              tools=[gmail_read_latest, send_email, forward_emails])
agent = secure_agent(agent, policy=policy)

Runner.run_sync(agent, "Please handle my latest unread email.")
# Output: IFCViolation: forbidden capabilities present in arguments
\end{lstlisting}

\paragraph{Flow.} Figure~\ref{fig:case-a} traces the labels at each
boundary. The injection never reaches the privileged tool because
the call carries the UNTRUSTED label inherited from
\texttt{gmail\_read\_latest}.

\begin{figure}[t]
  \centering
  \includegraphics[width=0.92\textwidth]{figures/Figure2}
  \caption{Scenario A. The email body is labelled UNTRUSTED at the
  gmail-read boundary. The label propagates through the planner's
  reasoning to the \texttt{forward\_emails} call. The IFC engine
  rejects the call because the tool requires control-flow integrity
  TRUSTED, blocking the attack by construction.}
  \label{fig:case-a}
\end{figure}


\subsection{Case B: PII leakage in a streaming voice agent}
\label{case:stream}

\paragraph{Persona.} Daniel is an ML engineer at a voice-AI startup.
The product streams customer-service responses token by token over a
phone call. Time-to-first-token must stay below 800 ms. He tried
Llama Guard 4 as an output guardrail but its 1.4 s full-response
latency was unusable, so he disabled it before launch.

\paragraph{Threat.} A user asks about refund processing. The model
hallucinates an example credit-card number that it learned from
training contamination. Without an output guardrail, the number
streams to the caller before any monitoring catches it.

\paragraph{Before.} Listing~\ref{lst:daniel-before} shows Daniel's
streaming loop after he disabled the heavyweight guardrail. There is
nothing between the model and the audio synthesizer.

\begin{lstlisting}[caption={Voice agent before AIsafePy (Scenario B).}, label={lst:daniel-before}]
async for token in model.stream(prompt):
    audio_synth.feed(token)
    # No guardrail in the loop: anything the model emits reaches
    # the caller within ~80 ms.
\end{lstlisting}

\paragraph{After.} Listing~\ref{lst:daniel-after} wraps the stream in
a \texttt{GuardPipeline} with a 50,000-char deterministic regex
guard at Tier 1, a small Prompt Guard 2 classifier at Tier 2, and
Llama Guard 4 at Tier 3 (which only runs on the 1-3\% of responses
that fall inside the uncertainty band). Tier 1 catches the
credit-card pattern within the first emitted chunk.

\begin{lstlisting}[caption={Voice agent with AIsafePy.stream (Scenario B).}, label={lst:daniel-after}]
from aisafepy.stream import GuardPipeline, RegexGuard, ClassifierGuard

pipeline = GuardPipeline(
    tier1=[RegexGuard.compile_pii(redact=False)],
    tier2=[ClassifierGuard.from_hf("meta-llama/Llama-Prompt-Guard-2-22M")],
    tier3=[ClassifierGuard.from_hf("meta-llama/Llama-Guard-4")],
    budget_ms_p95=80,
)

async for piece in pipeline.guard_stream(model.stream(prompt), prompt=prompt):
    if hasattr(piece, "action"):
        # GuardDecision returned: block downstream synthesis.
        log_to_observability(piece)
        audio_synth.fallback("I cannot share that information.")
        break
    audio_synth.feed(piece)
\end{lstlisting}

\paragraph{Flow.} Figure~\ref{fig:case-b} shows the cascade in
action. Tier 1 regex catches the credit-card span at token 4 and
the cascade aborts; the audio synthesizer is told to play a
fallback instead of the emitted tokens. End-to-end added latency
on the blocked path is roughly 5 ms.

\begin{figure}[t]
  \centering
  \includegraphics[width=0.92\textwidth]{figures/Figure3}
  \caption{Scenario B. Streaming tokens flow into the Tier 1 regex
  guard. When a credit-card pattern is matched mid-stream, the
  pipeline emits a terminal \texttt{GuardDecision} and the
  downstream audio synthesizer receives a fallback string instead
  of the leaked tokens.}
  \label{fig:case-b}
\end{figure}


\subsection{Case C: PyRIT failures into runtime guards}
\label{case:adapt}

\paragraph{Persona.} Priya runs the security team for a B2B SaaS
that recently shipped an AI assistant on top of its docs. She runs
PyRIT \citep{pyrit} every two weeks. The last campaign produced
about 5,000 attempts and 312 policy violations. Translating those
into Guardrails-AI validators takes her roughly two engineer-days
per sprint and lags real attacker mutation by ten days.

\paragraph{Threat.} A red-team campaign finds 47 successful
base64-encoded jailbreak payloads. The window between
discovery and deployment of a runtime guard is what attackers
exploit.

\paragraph{Before.} Listing~\ref{lst:priya-before} sketches Priya's
current workflow: read JSONL, hand-craft regex and validators, hope
nothing is missed.

\begin{lstlisting}[caption={Manual red-team to guardrail workflow (Scenario C).}, label={lst:priya-before}]
# 1. Read 5000 lines of pyrit_memory.duckdb manually.
# 2. Sort by scorer label, cluster mentally.
# 3. For each cluster, hand-write a Guardrails-AI validator:
@register_validator(name="base64-jailbreak", data_type="string")
class Base64Jailbreak(Validator):
    def validate(self, value, metadata): ...
# 4. Deploy. Wait for the next red-team campaign. Repeat.
\end{lstlisting}

\paragraph{After.} Listing~\ref{lst:priya-after} replaces that
workflow with a single \texttt{compiler.compile()} call. The output
is a typed \texttt{CompilationReport} containing four artifact
modalities per cluster: a distilled classifier, synthesized regex
patterns, a Cedar policy rule, and a deliberative case appended to
the system prompt.

\begin{lstlisting}[caption={Eval-to-guard compilation with AIsafePy.adapt (Scenario C).}, label={lst:priya-after}]
from aisafepy.adapt import PyRITSource, GuardCompiler, Target, promote
from aisafepy.stream import GuardPipeline

source = PyRITSource(memory_db="pyrit_memory.duckdb", only_violations=True)

compiler = GuardCompiler(
    source=source,
    targets=[
        Target.distill_classifier(base="meta-llama/Llama-Prompt-Guard-2-22M"),
        Target.synthesize_regex(min_precision=0.99),
        Target.policy_rule(dsl="cedar"),
        Target.deliberative_case(policy="policies/company_safety.md"),
    ],
    min_attack_success_rate=0.05,
)
report = compiler.compile()
promote(report, to=GuardPipeline.from_yaml("guards.yaml"),
        canary_traffic_pct=1.0, fp_budget=0.005)
# Projected cluster ASR on holdout: 71% -> 6%
\end{lstlisting}

\paragraph{Flow.} Figure~\ref{fig:case-c} shows the adapt pipeline:
ingestion, clustering, per-target compilation, and canary deploy
with a false-positive budget.

\begin{figure}[t]
  \centering
  \includegraphics[width=0.92\textwidth]{figures/Figure4}
  \caption{Scenario C. The adapt compiler ingests PyRIT failure
  records, embeds them with sentence-transformers, clusters with
  HDBSCAN, and emits one or more artifact types per cluster
  (distilled classifier, regex, Cedar policy, steering vector,
  deliberative case). The canary stage shadows new artifacts under
  a false-positive budget before promoting them to production.}
  \label{fig:case-c}
\end{figure}
